\documentclass[12pt,a4paper]{article}
\usepackage[french]{babel}
\usepackage[T1]{fontenc}
\usepackage{amsmath, amssymb}
\usepackage{geometry}
\geometry{top=1cm, bottom=2cm, left=1cm, right=1cm}
\usepackage{tcolorbox}
\usepackage{listings}
\usepackage{xcolor}
%\usepackage{multicol}

\definecolor{codegray}{rgb}{0.95,0.95,0.95}
\lstset{
    language=Python,
    backgroundcolor=\color{codegray},
    basicstyle=\ttfamily\small,
    keywordstyle=\color{blue},
    commentstyle=\color{gray},
    stringstyle=\color{red},
    frame=single,
    breaklines=true
}

\begin{document}

\centering\textbf{Devoir Surveillé d’Algorithmique - Python : Affectation et condition \texttt{if}}\\
\vspace{0.5cm}

\noindent\textbf{Nom :} \underline{\hspace{5cm}} \hfill \textbf{Prénom :} \underline{\hspace{5cm}} \\
\textbf{Classe :} \underline{\hspace{5cm}}\\
\vspace{1cm}

%\begin{multicols}{2}
\section*{Sujet A}

\begin{enumerate}
    \item Que va afficher le programme suivant après exécution ? Justifie ta réponse.

\begin{lstlisting}
note = 14
moyenne = 10

if note >= moyenne:
    print("Admis")
    if note >= 16:
        print("Mention tres bien")
else:
    print("Refuse")
\end{lstlisting}

\begin{tcolorbox}[colframe=black, colback=white, width=\linewidth]
\vspace{3.5cm}
\end{tcolorbox}

    \item Écris un programme en Python qui :
    \begin{itemize}
        \item crée deux variables \texttt{a} et \texttt{b}, avec les valeurs de ton choix ;
        \item affiche "\texttt{La somme est paire}" si la somme est paire ;
        \item sinon affiche "\texttt{La somme est impaire}".
    \end{itemize}

\begin{tcolorbox}[colframe=black, colback=white, width=\linewidth]
\vspace{5cm}
\end{tcolorbox}

\end{enumerate}

\vspace{0.3cm}
\textbf{Rappel :} Pour savoir si un nombre est pair, on utilise :\\
\texttt{if nombre \% 2 == 0:}

%\columnbreak
\newpage

\section*{Sujet B}

\begin{enumerate}
    \item Que va afficher le programme suivant après exécution ? Justifie ta réponse.

\begin{lstlisting}
note = 9
moyenne = 10

if note >= moyenne:
    print("Admis")
    if note >= 16:
        print("Mention tres bien")
else:
    print("Refuse")
\end{lstlisting}

\begin{tcolorbox}[colframe=black, colback=white, width=\linewidth]
\vspace{3.5cm}
\end{tcolorbox}

    \item Écris un programme en Python qui :
    \begin{itemize}
        \item crée deux variables \texttt{x} et \texttt{y}, avec les valeurs de ton choix ;
        \item affiche "\texttt{La somme est paire}" si la somme est paire ;
        \item sinon affiche "\texttt{La somme est impaire}".
    \end{itemize}

\begin{tcolorbox}[colframe=black, colback=white, width=\linewidth]
\vspace{5cm}
\end{tcolorbox}

\end{enumerate}

\vspace{0.3cm}
\textbf{Rappel :} Pour savoir si un nombre est pair, on utilise :\\
\texttt{if nombre \% 2 == 0:}

%\end{multicols}

\end{document}
